%% ============================================================
%% CHAPTER DRAFT — VP2: Integration and Visualization of 3D Data
%%                 for the Ballenberg Open-Air Museum
%%
%% Usage: \input{chapters_draft} in pointclouds.tex after the
%%        State of the Art chapter.
%% ============================================================


% ============================================================
\chapter{System Architecture and Implementation}
\label{ch:implementation}
% ============================================================

This chapter describes the platform that was designed and built during the
implementation workpackage (calendar weeks 13--17). The system addresses the
pipeline and visualization challenges identified in the State of the Art review:
it converts heterogeneous survey data to 3D Tiles for web delivery (Research
Question 3), provides a web-based CesiumJS viewer with measurement, annotation,
and analysis tools suited to the target user groups (Research Question 2), and
introduces a structured dataset registry with a minimum metadata schema as a
first step toward a scalable archive (Research Question 1).

All source code is maintained in the GitHub repository
\texttt{joelbeachman/pointcloud-platform} and runs inside a Docker container
for reproducible deployment.

% ------------------------------------------------------------
\section{Platform Architecture}
\label{sec:arch}
% ------------------------------------------------------------

The system follows a lightweight three-tier architecture
(Figure~\ref{fig:architecture}): a Node.js/Express HTTP server handles the REST
API and serves static assets; a JSON-based dataset registry on disk stores
metadata for all registered datasets; and a browser-based frontend delivers
the interactive viewers.

\begin{figure}[h]
\centering
% Placeholder: replace with actual architecture diagram
\fbox{\parbox{0.7\textwidth}{\centering\vspace{2cm}%
  [Architecture diagram: Browser → Express server → Data directory /
  datasets.json → 3D Tiles / panoramas / splats]\vspace{2cm}}}
\caption{High-level platform architecture. All components run in a single
  Docker container; the data volume is mounted externally.}
\label{fig:architecture}
\end{figure}

\subsection{Server and REST API}

The server (\texttt{server.js}) is built on Express and exposes the following
principal endpoints:

\begin{itemize}
  \item \texttt{GET /api/datasets} — returns all registered datasets.
  \item \texttt{POST /api/datasets} — registers a new dataset entry.
  \item \texttt{PATCH /api/datasets/:id} — updates metadata, including the
    \texttt{modelMatrix} that aligns the tileset to the LV95 coordinate
    reference system.
  \item \texttt{DELETE /api/datasets/:id} — removes a dataset from the
    registry.
  \item \texttt{POST /api/helmert} — computes a six- or seven-parameter
    Helmert transformation from ground control point pairs and returns the
    rotation matrix, scale, translation, per-point residuals, and RMS error.
  \item \texttt{POST /api/profile} — elevation profile query: accepts a polyline
    drawn on the scene and a half-width, reads the relevant \texttt{.pnts} tiles
    from disk, and returns projected $(d,\,z)$ point data for display
    (Section~\ref{sec:profile}).
  \item \texttt{GET /cesium-proxy/*} — transparently proxies the CesiumJS
    library and the \texttt{decodeSpz.js} Web Worker from the CDN, patching
    both with safe typed-array helpers. The patch replaces calls to
    \texttt{.typedArray.slice()} and \texttt{.value.slice()} that occur in
    the Gaussian splat decoder path, where the Emscripten WASM module returns
    vector objects that do not implement the standard \texttt{.slice()} method,
    causing a runtime error when loading splat assets. All other workers and
    static assets (\texttt{Assets/}, \texttt{ThirdParty/}) are redirected to
    the CDN unchanged.
\end{itemize}

\subsection{Dataset Registry}

All datasets are described in a single file \texttt{data/datasets.json}. Each
entry contains at minimum the fields listed in Table~\ref{tab:metadata_min},
which aligns with the minimum metadata schema proposed in
Section~\ref{sec:structuring} of the State of the Art review. Optional fields
include \texttt{modelMatrix} (a 4\,$\times$\,4 column-major LV95\,$\to$\,ECEF
transformation matrix), \texttt{helmertLV95} (the seven Helmert parameters
used to derive it), \texttt{panoramas} (scan station positions and paths for
equirectangular images), and a set of Ballenberg-specific provenance fields
(\texttt{buildingId}, \texttt{captureDate}, \texttt{captureMethod},
\texttt{scannerModel}).

\begin{table}[h]
\centering
\caption{Minimum required metadata fields per dataset entry.}
\label{tab:metadata_min}
\begin{tabular}{lll}
\toprule
\textbf{Field} & \textbf{Type} & \textbf{Description} \\
\midrule
\texttt{id}         & string    & Unique identifier (slug) \\
\texttt{name}       & string    & Human-readable label \\
\texttt{type}       & enum      & \texttt{cesium | cesium-splat | splat | e57 | pointcloud | potree} \\
\texttt{path}       & string    & Server-relative path to the primary asset \\
\texttt{createdAt}  & ISO 8601  & Registration timestamp \\
\bottomrule
\end{tabular}
\end{table}

\subsection{Coordinate Reference System}

All geospatial datasets in the Ballenberg project are captured in the Swiss
national reference frame LV95 (EPSG:2056), which uses East ($E$), North ($N$),
and ellipsoidal height ($H$) coordinates in metres with values on the order of
$2.6 \times 10^6$\,m East and $1.2 \times 10^6$\,m North. CesiumJS, in
contrast, places all geometry in ECEF (Earth-Centered, Earth-Fixed) coordinates.

The model matrix that transforms LV95 absolute coordinates into ECEF is
constructed via:
\begin{equation}
  \mathbf{M}_{\text{LV95}} = \mathbf{M}_{\text{ENU}}(E_0, N_0, H_0)
    \cdot \mathbf{T}(-E_0, -N_0, -H_0),
\end{equation}
where $(E_0, N_0, H_0)$ is the bounding-sphere centre of the tileset in
LV95 coordinates, converted to WGS\,84 via the approximate
LV95\,$\to$\,WGS\,84 formula implemented in \texttt{lv95ToWgs84()}, $\mathbf{M}_{\text{ENU}}$ is the
east-north-up frame at that ECEF point, and $\mathbf{T}$ is a translation
that brings the origin of the LV95 frame to the centroid. The effect is that
the tileset's bounding sphere sits at the correct geographic position without
any offset, and plane or shader computations operating in tileset-local
coordinates directly work in metric East/North/Up units centred on the
dataset.

% ------------------------------------------------------------
\section{Data Processing Pipeline}
\label{sec:pipeline}
% ------------------------------------------------------------

The processing script \texttt{scripts/process.py} implements a universal
converter that auto-detects the format of a given input file and produces
the appropriate streaming asset together with a registry entry. The following
input modalities are supported:

\begin{description}
  \item[Point clouds (LAS/LAZ, E57, XYZ/TXT/PTS, PTX, PCD, PLY)] — converted
    to a single-tile 3D Tiles 1.0 \texttt{.pnts} asset with an accompanying
    \texttt{tileset.json}. XYZ colour attributes are preserved; the floor of
    the bounding-box minimum coordinate is stored as \texttt{RTC\_CENTER},
    keeping all float32 position offsets small and positive. A LV95 model
    matrix is computed and stored.
  \item[Meshes (OBJ, STL, GLB, GLTF, mesh PLY)] — converted to a
    \texttt{.b3dm}-based 3D Tiles scene.
  \item[Gaussian splats (.splat, 3DGS PLY)] — \texttt{.splat} inputs are
    copied directly; 3DGS PLY inputs are converted by decoding the DC
    spherical harmonics coefficients ($f\_dc\_0/1/2$) to RGB via a sigmoid
    activation, normalising the quaternion rotation, and packing all
    attributes into the compact 32-byte-per-Gaussian layout expected by
    \texttt{@mkkellogg/gaussian-splats-3d}.
\end{description}

\subsection{Haus Eggiwil Dataset}

The primary real-world dataset used during the implementation phase is the
terrestrial laser scan of Haus Eggiwil (building code used internally:
\texttt{haus-eggiwil}), captured with a Leica scanner and stored in a LAS file
\texttt{351\_Haus-Eggiwil.las}. After processing, the resulting 3D Tiles asset
contains 5\,023\,669 points covering the building exterior and immediate
surroundings. In addition, 185 panoramic scan positions were extracted from the
E57 archive and stored as equirectangular JPEGs in
\texttt{data/panoramas/haus-eggiwil/}, with a \texttt{metadata.json} file
encoding LV95 coordinates, scanner orientation, and north-offset angles.

The north-offset angles — required to orient Pannellum's equirectangular
viewer to true north — were derived directly from the quaternion columns of
the \texttt{image\_poses.csv} file using
\begin{equation}
  \alpha_N = \operatorname{atan2}\!\left(1 - 2(q_y^2 + q_z^2),\;
    2(q_x q_y + q_w q_z)\right),
\end{equation}
which avoids the ambiguity introduced by Euler ZYX decomposition near
$|R| \approx 180°$ (gimbal-lock variant).

% ------------------------------------------------------------
\section{Web Visualization Layer}
\label{sec:viewer}
% ------------------------------------------------------------

The primary viewer is built on CesiumJS 1.140 (Apache 2.0), served through the
local proxy that caches and patches the library. The viewer window fills the
area between two 220\,px sidebars: a left panel for layer management and tools,
and a right panel for measurement results.

\subsection{Multi-Layer Management}

Any number of datasets can be added as layers via the portal dashboard or the
in-viewer ``Add Layer'' button. Each layer maintains its own CesiumJS
\texttt{Cesium3DTileset} instance. The viewer supports three layer types:

\begin{itemize}
  \item \textbf{Cesium (3D Tiles)} — loaded via \texttt{Cesium3DTileset} and
    positioned via the LV95 model matrix.
  \item \textbf{Gaussian Splat} — rendered by a Three.js canvas overlaid on
    the Cesium canvas; the viewer cameras are synchronized every frame by
    decomposing the Cesium camera ECEF matrix into the ENU frame at the
    dataset centroid and transferring the resulting translation and quaternion
    to the Three.js camera.
  \item \textbf{Panorama} — scan station billboards added to the Cesium scene;
    clicking a billboard opens the equirectangular overlay.
\end{itemize}

\subsection{Coordinate Alignment: Helmert Registration}

Datasets that arrive in a scanner-local coordinate system rather than LV95
absolute can be aligned using the in-viewer GCP registration tool. The user
clicks corresponding points in the Cesium scene (ECEF, converted internally to
tileset-local) and enters the known LV95 coordinates. A six- or seven-parameter Helmert transformation (rotation and translation,
optionally with uniform scale) is computed via SVD of the centred point pairs
by the server-side \texttt{POST /api/helmert} endpoint; the resulting
rotation matrix $\mathbf{R}$, scale $s$, and translation $\mathbf{t}$ are
returned to the client, which assembles the $4 \times 4$ model matrix
$\mathbf{M} = \mathbf{M}_{\text{LV95}} \cdot \mathbf{M}_{\text{Helmert}}$
and stores it persistently in the dataset registry.

% ------------------------------------------------------------
\section{Measurement Tools}
\label{sec:measurement}
% ------------------------------------------------------------

Five measurement tools are available in the left panel and operate on pick
positions obtained from \texttt{scene.pickPosition()} with a bounding-sphere
fallback for Gaussian splat layers (which do not populate the depth buffer):

\begin{enumerate}
  \item \textbf{Distance} — 3D straight-line distance between two clicked
    points, displayed as a labelled polyline entity.
  \item \textbf{Horizontal distance} — horizontal plan distance (XY component
    only in tileset-local space).
  \item \textbf{Vertical distance} — signed height difference ($\Delta Z$).
  \item \textbf{Area} — polygon area computed via the shoelace formula on
    points projected into the local XY plane; right-click or double-click
    finishes the polygon.
  \item \textbf{Coords} — displays LV95 coordinates and ellipsoidal height for
    a clicked point; the ECEF position is converted back through
    $\mathbf{M}_{\text{LV95}}^{-1}$.
\end{enumerate}

All committed measurements are stored in an in-memory array and rendered both
in the 3D Cesium scene (polyline/polygon entities) and projected onto the
panorama canvas whenever the overlay is open.

% ------------------------------------------------------------
\section{Integrated Panorama Overlay}
\label{sec:pano}
% ------------------------------------------------------------

The panorama viewer is embedded as an overlay covering the central viewer
canvas (leaving both sidebars visible) and powered by Pannellum~2.5.6
(MIT licence). Clicking a camera-marker billboard in the Cesium scene opens
the corresponding equirectangular image at the correct initial heading.

\subsection{Spatial Navigation}

Hotspot icons appear at the projected positions of nearby scan stations
(horizontal distance $\leq 10$\,m). The bearing to each neighbour is computed
as:
\begin{equation}
  \beta = \operatorname{atan2}(\Delta E,\, \Delta N) - \alpha_{N,\,\text{src}},
\end{equation}
where $\alpha_{N,\,\text{src}}$ is the north-offset of the current station.
Icon size scales with inverse distance (56\,px at 0.5\,m to 24\,px at
$\geq 6.25$\,m). On navigation, the heading offset relative to the source
view direction is preserved, so the user arrives at the next station looking
in the same geographic direction.

\subsection{Panorama Measurement Integration}

When a measurement tool is active, clicking inside the panorama casts a ray
from the scanner position in the direction given by the inverse perspective
projection of the click:

\begin{equation}
  \hat{\mathbf{w}} = \frac{\Delta x}{f}\hat{\mathbf{r}}
    - \frac{\Delta y}{f}\hat{\mathbf{u}} + \hat{\mathbf{f}},
  \quad
  f = \frac{W/2}{\tan(\text{HFOV}/2)},
\end{equation}
where $\hat{\mathbf{r}}$, $\hat{\mathbf{u}}$, $\hat{\mathbf{f}}$ are the
right, up, and forward unit vectors of the current Pannellum view in world
space, and $(\Delta x, \Delta y)$ is the click offset from the image centre.
The world-space ray is then intersected with the Cesium scene via
\texttt{scene.pickFromRay()}, or falls back to a floor-plane intersection at
a configurable height offset when the ray misses the point cloud geometry.

Committed measurement markers are reprojected onto the panorama canvas each
animation frame so they remain visible regardless of whether the points were
placed in the 3D view or the panorama overlay.

% ------------------------------------------------------------
\section{Clip Box Tool}
\label{sec:clipbox}
% ------------------------------------------------------------

The clip box tool allows any number of axis-aligned boxes — optionally rotated
— to be defined and applied simultaneously as spatial filters on all loaded
tilesets.

\subsection{Box Definition and Rotation}

Each box is parameterized by six face offsets ($x_{\min}$, $x_{\max}$,
$y_{\min}$, $y_{\max}$, $z_{\min}$, $z_{\max}$) in ENU-relative metres from
the dataset centroid, plus three rotation angles $(\theta_x, \theta_y,
\theta_z)$ in degrees. The combined rotation matrix is:
\begin{equation}
  \mathbf{R} = \mathbf{R}_z(\theta_z)\,\mathbf{R}_y(\theta_y)\,
    \mathbf{R}_x(\theta_x) \quad \text{(row-major)},
\end{equation}
and the six face normals in tileset-local space are the column vectors of
$\mathbf{R}$, scaled and offset to the face positions.

\subsection{Clipping Modes and Boolean Combinations}

Each box can be set to one of three modes:
\begin{description}
  \item[Inside] — retains only points within the box. The six rotated plane
    normals point outward; \texttt{unionClippingRegions = true} clips anything
    outside any plane (intersection of the six half-spaces).
  \item[Outside] — removes points within the box (subtraction). Plane normals
    point inward; \texttt{unionClippingRegions = false}.
  \item[None] — the box acts as a visual reference only (wireframe displayed,
    no clipping planes applied).
\end{description}

When only inside boxes or only outside boxes are active, the standard
\texttt{ClippingPlaneCollection} API is used, which operates at tile-level
culling granularity (fast). When both inside and outside boxes are
simultaneously active — a boolean combination that cannot be expressed as a
flat list of planes — a \texttt{Cesium.CustomShader} is generated at runtime.
The fragment shader receives each point's position in eye space from
\texttt{czm\_modelView}, transforms it to ENU via a per-frame uniform
\texttt{u\_eyeToEnu}:
\begin{equation}
  \mathbf{M}_{\text{eye}\to\text{ENU}} = \mathbf{M}_{\text{ECEF}\to\text{ENU}}
    \cdot \mathbf{M}_{\text{eye}}^{-1},
\end{equation}
and evaluates all box conditions per fragment with \texttt{discard}
statements. This avoids the float32 ECEF precision errors ($\approx 0.5$\,m at
Ballenberg latitudes) by keeping numerical operations in camera-relative eye
space, where the translation to the scene origin is already encoded in the
view matrix.

\subsection{Interactive 3D Handles}

Six face handles (coloured red/green/blue by axis) and one centre handle are
placed at the face and centroid positions of the active box. Dragging a face
handle constrains movement along the corresponding local axis. Dragging the
centre handle translates the entire box. Three rotation rings (one per axis)
allow free angular adjustment by projecting the mouse arc onto the screen
plane at the box centre. All geometry updates — wireframe, clipping planes,
and editor inputs — are applied continuously during drag.

% ------------------------------------------------------------
\section{Elevation Profile}
\label{sec:profile}
% ------------------------------------------------------------

The elevation profile tool (Ger.\ \textit{Höhenprofil}) provides a
functionality similar to Potree's built-in cross-section profile. It is
motivated by Research Questions 2 and 3: it directly addresses the
conservation and survey use case of extracting a vertical section through a
point cloud along an arbitrary path.

\subsection{Workflow}

The user activates the Profile tool in the left panel and clicks waypoints on
the Cesium scene to define a multi-segment profile line (right-click finishes
the line). Each waypoint position is converted from ECEF to tileset-local
space via $\mathbf{M}_{\text{LV95}}^{-1}$ before being sent to the server.

\subsection{Server-Side Point Extraction}

The \texttt{POST /api/profile} endpoint:
\begin{enumerate}
  \item Reads the dataset's \texttt{tileset.json} and walks the tile tree,
    accumulating the column-major $4 \times 4$ tile transform matrices at
    each node.
  \item For each leaf \texttt{.pnts} tile, reads the 28-byte header and
    the feature table JSON to extract \texttt{POINTS\_LENGTH},
    \texttt{RTC\_CENTER}, and the byte offsets for \texttt{POSITION}
    (float32 XYZ) and \texttt{RGB} (uint8).
  \item Streams positions in blocks of 262\,144 points ($\approx 3$\,MB per
    read), applies the cumulative tile transform, and for each point computes
    the perpendicular distance to the nearest profile segment:
    \begin{equation}
      d_\perp = \min_k \left\|
        (p_{xy} - q_k)\,-\,
        \frac{(p_{xy}-q_k)\cdot\hat{s}_k}{\|s_k\|}\hat{s}_k
      \right\|,
    \end{equation}
    where $q_k$ is the start of segment $k$ and $\hat{s}_k$ its unit direction.
  \item Points with $d_\perp \leq w/2$ (half the user-specified slab width)
    are retained; their projected distance along the polyline $d$ and
    elevation $z$ are recorded together with the RGB colour.
  \item Up to 150\,000 output points are returned as JSON, sorted by $d$.
\end{enumerate}

\subsection{Profile Panel}

A floating panel at the bottom of the viewer renders the profile as a 2D
scatter plot (elevation vs.\ distance along line) on an HTML5 canvas. Points
are drawn as $2 \times 2$\,px dots using the \texttt{ImageData} API for
efficient rendering at 150\,000 points. If the dataset has RGB attributes, the
original point colours are used; otherwise, points are coloured by a
Viridis-like gradient mapping elevation to colour. Axis labels, a hover
tooltip showing $(d,\,z)$, and a width control are provided. The panel can be
dismissed with Escape.

\medskip

Table~\ref{tab:tools} provides a summary of all implemented interactive tools
and their primary use cases for the target user groups identified in
Section~\ref{sec:introduction}.

\begin{table}[h]
\centering
\caption{Summary of implemented interactive tools and target user groups.}
\label{tab:tools}
\begin{tabular}{llll}
\toprule
\textbf{Tool} & \textbf{Panel} & \textbf{Primary use case} & \textbf{User group} \\
\midrule
Distance, Horizontal, Vertical & Left & Dimension checking & Conservation staff \\
Area                           & Left & Floor area estimation & Conservation staff \\
Coordinates                    & Left & Position inspection & Survey team \\
Elevation Profile              & Left & Cross-section analysis & Survey / conservation \\
Clip Box (Inside/Outside)      & Left & Selective visibility & All \\
Panorama Overlay               & Toolbar & Contextual imagery & All \\
Panorama Measurement           & Panorama & In-photo measurement & Survey team \\
GCP Registration (Helmert)     & Right & Coordinate alignment & Survey team \\
Scan Station Navigation        & Scene  & Site exploration & All \\
\bottomrule
\end{tabular}
\end{table}


% ============================================================
\chapter{Open Tasks and Future Work}
\label{ch:open_tasks}
% ============================================================

The following tasks remain open before the platform can be considered complete
for the purposes of this project. They are grouped by priority.

% ------------------------------------------------------------
\section{High Priority}
\label{sec:tasks_high}
% ------------------------------------------------------------

\subsection{GCP Registration: Missing \texttt{helmert.py}}

The in-viewer GCP georeferencing tool sends point-pair data to the
\texttt{POST /api/helmert} server endpoint, which is configured to execute
\texttt{scripts/helmert.py}. This script does not currently exist in the
repository, so any attempt to compute a Helmert transformation from the
viewer will return a \texttt{500} error. The SVD algorithm is already
implemented as a JavaScript function \texttt{solveHelmert()} in the viewer
source (kept as a fallback reference) and can be ported to a short Python
script using NumPy. This script must be created before the GCP registration
feature can be considered functional.

\subsection{Panorama Image Verification}

The equirectangular JPEGs stored in \texttt{data/panoramas/haus-eggiwil/}
have not yet been verified as geometrically correct equirectangular
projections. Leica scanners store panoramic imagery in a scanner-specific
format that may require resampling to true equirectangular before Pannellum
can display them without distortion. A systematic check — and, where
necessary, a reprojection step using the scanner's calibration parameters
— is required before the panorama tool can be considered validated on real data.

\subsection{North-Offset Recovery Script Update}

The script \texttt{scripts/restore\_eggiswil.py}, which documents and
automates the restoration of panorama metadata from backup, currently derives
the north-offset angle from the \texttt{rotZ\_deg} column of
\texttt{image\_poses.csv}. As noted in Section~\ref{sec:pipeline}, the
correct formula uses the quaternion columns directly. The script must be
updated to use Equation~(2) to ensure north-offsets can be recovered
consistently after any future data loss.

\subsection{Documentation Folder Restoration}

Thesis source assets — figures, compiled PDF, and ancillary documentation
— were lost during a \texttt{git reset --hard} operation in calendar
week 15. These must be recovered from external backup before the final report
can be compiled. A \texttt{.gitignore} rule and a data safety protocol have
since been established to prevent recurrence.

\subsection{End-to-End Validation and v0.1.0 Tag}

Before tagging a stable release, the full pipeline from raw LAS input to
interactive CesiumJS viewer should be run end-to-end on the Haus Eggiwil
dataset with a fresh environment to verify reproducibility. This includes
verifying that the LV95 model matrix, panorama positions, north-offsets, and
measurement tools all produce geometrically consistent results at the same
scan position.

% ------------------------------------------------------------
\section{Medium Priority}
\label{sec:tasks_medium}
% ------------------------------------------------------------

\subsection{COPC Streaming Support}

The current processing pipeline generates a single-tile \texttt{.pnts} asset
for each point cloud dataset. For large datasets such as a multi-building
Ballenberg scan (tens of millions of points), single-tile loading creates
significant latency and limits progressive rendering. Converting LAS files to
Cloud Optimized Point Cloud (COPC) format, as discussed in
Section~\ref{sec:formats}, would enable streaming from the existing server
without structural changes. COPC is directly supported by recent versions of
PotreeConverter and can be served as a static file from any HTTP server.

\subsection{Metadata Schema Validation}

The minimum metadata schema defined in Table~\ref{tab:metadata_min} has been
implemented in the dataset registry but has not been validated against the
actual documentation and conservation workflows of the Ballenberg museum. A
structured review with domain experts — identifying which additional fields
(capture method, scanner model, coordinate accuracy, conservation status) are
necessary for operational use — is required before the schema can be
considered stable. The schema should ultimately be alignable with
CIDOC-CRM~/ CRMdig, as identified in Gap~2 of the synthesis
(Section~\ref{sec:synthesis}).

\subsection{Semantic Annotations Layer}

A clickable annotation layer — allowing conservators to attach free-text notes,
image links, or structured condition reports to spatial regions in the Cesium
viewer — would directly address Research Question~4 (semantic metadata linked
to 3D visualization). A minimal viable implementation would store annotations
as GeoJSON-compatible objects in the dataset registry, render them as
\texttt{BillboardCollection} entities in CesiumJS, and display a detail panel
on click. Full CIDOC-CRM property sets would be a subsequent extension.

\subsection{Mobile Optimization Audit}

The current viewer has not been tested on mobile browsers. As noted in
Section~\ref{sec:visualization}, mobile performance for WebGL point cloud
viewers varies strongly by device. CesiumJS 1.140 has better mobile
compatibility than Potree~1.8, but the current viewer uses several features
(Gaussian splat Three.js canvas overlay, WebGL depth buffer read for pick
positions) that may degrade on mobile GPUs. A systematic audit and
optimization pass — reducing shader complexity, disabling the splat overlay
on low-memory devices, and testing on representative Android and iOS hardware
— is required before the platform can be offered to museum visitors.

% ------------------------------------------------------------
\section{Long-Term / Future Research}
\label{sec:tasks_future}
% ------------------------------------------------------------

\subsection{Multi-Building Archive Structure}

The current dataset registry is flat: each entry represents one dataset
regardless of its relationship to a building, site, or documentation campaign.
For a full Ballenberg deployment covering more than 100 buildings, a
hierarchical structure — grouping datasets by building ID, campaign, and
processing level — is required. This connects directly to the file naming and
folder structure conventions discussed in Section~\ref{sec:structuring}. A
migration of the current flat JSON registry to a more structured schema (or a
lightweight SQLite database) should be planned once the number of registered
datasets exceeds a practical threshold for manual management.

\subsection{Potree-Next Integration Pathway}

As discussed in Section~\ref{subsec:potree}, the Potree-Next rewrite is
progressing toward WebGPU-based rendering with compute shaders and explicit
mobile support. Once Potree-Next reaches a stable release with feature parity
to Potree~1.8, the point cloud viewer layer of this platform should be
re-evaluated for migration. The 3D Tiles assets and metadata already generated
are format-compatible, so the migration effort is concentrated in the viewer
layer rather than the data layer.

\subsection{CIDOC-CRM Metadata Export}

The dataset-level metadata currently stored in \texttt{datasets.json} contains
the fields required for a CIDOC-CRM serialization (capture method, date,
instrument, coordinate reference system, processing status). A machine-readable
export in JSON-LD conforming to the CRMdig extension would allow the Ballenberg
inventory to be linked to national and European heritage databases. This
addresses Gap~4 from the synthesis and is a longer-term research deliverable
beyond the scope of VP2.

\subsection{User Study}

The three target user groups hypothesized in Section~\ref{sec:introduction} —
conservation staff, museum visitors, and student workshop participants — have
not yet been evaluated. A structured user study, using the implemented platform
as a prototype, would validate or revise the interaction requirements, identify
missing functionality, and provide the evidence base for Research
Question~2. This is the primary task of the evaluation workpackage scheduled
for calendar weeks 17--20.
